\chapter{Conclusion}

\section{Key Findings}
A Tunisian entrant cannot win by copying an international AMR and undercutting on price -- components are imported, cost structure is not competitive. Advantages lie in modular design, local manufacturing/support, EU proximity, and pricing tailored to the local market.

Sensor/inspection is the lowest-friction entry, requiring no facility changes. Manufacturing intralogistics (for automotive exporters) and universities are the strongest initial markets, with manufacturing providing a funded mid-term expansion path.
\section{Recommendation}
\begin{itemize}
    \item \textbf{Immediate}: industrial inspection + university pilots in Tunisia. ROI works without matching international prices -- inspection against avoided cost, universities against imported research robot cost. Fast decisions; build with existing ROS2/Nav2 stack, simulation, modular chassis.
    \item \textbf{Years 2--3}: export-oriented automotive/electronics intralogistics (arm + base) and 3PL/internal logistics (base + compartment), offered as a service to bypass capex.
    \item \textbf{Later}: GCC export (compartment/sensor configurations), private healthcare, and hospitality via international chains already in Tunisia.
    \item \textbf{Defer}: shelf-access module to a partner-funded phase. \textbf{Drop}: standalone security and in-store retail.
\end{itemize}

\section{Cross-Cutting Conditions}

Across all sectors:
\begin{enumerate}
    \item \textbf{Service-based pricing} (RaaS, usage-based, shared) \cite{builtin_raas}.
    \item \textbf{Hybrid build}: import only critical components; fabricate structure/modules locally.
    \item \textbf{French/Arabic interfaces} and visible local service.
    \item \textbf{Use Startup Act and EU-co-funded pilots} for financing.
\end{enumerate}
